% ============================================================
%  CHAPTER 1: INTRODUCTION
% ============================================================
\chapter{Introduction}
\label{ch:intro}

Since its creation, the internet has grown into a critical piece of global infrastructure that has enabled
unprecedented communication, commerce, and innovation, connecting over 29 billion devices worldwide~\cite{cisco2020air}.
This growth has been accompanied by a parallel rise in cyber crime, with global damages estimated at \$10.5 trillion annually by
2025~\cite{cybersecventures2025}.
Software vulnerabilities are a primary enabler of this damage: they provide the footholds through which attackers breach systems,
escalate privileges, and exfiltrate data.
The National Vulnerability Database recorded over 40,000 new CVEs in 2024 alone, a 38\% increase over the prior year~\cite{nvd2024stats},
and this count only reflects vulnerabilities that have been publicly reported.

The security community has responded with automated techniques for finding vulnerabilities at scale.
Static analyzers scan entire codebases for dangerous patterns.
Coverage-guided fuzzers generate millions of test inputs per second.
Symbolic execution engines reason about program paths.
These tools have stretched the volume of vulnerability detection far beyond what manual auditing can sustain,
and they continue to improve.
Yet a persistent gap separates \emph{detecting} a potential vulnerability from \emph{confirming} that it is exploitable.
Static analyzers produce candidates that may be infeasible.
Fuzzers produce crashes that may not be security-relevant.
Symbolic execution engines often cannot scale to whole programs.
Each technique, in isolation, yields partial evidence.

Vulnerability researchers generate evidence of vulnerabilities not through any single technique, but through a combination of approaches.
When an analyst examines a static analysis report, they do not simply check whether a data flow
exists from source to sink; they evaluate what happened to the value along the way, whether
the path is reachable under realistic conditions, and how findings from one technique corroborate
those of another.
These judgments, about value transformation, path feasibility, and cross-technique corroboration, are the
difference between a list of candidates and a feasible vulnerability.


\section{Firmware Vulnerability Discovery}

The Internet of Things has placed billions of devices into homes, offices, and critical infrastructure, with projections exceeding 30 billion connected
devices by the end of the decade~\cite{statista2024iot}.
These devices run firmware that is frequently complex, feature-rich, and developed under resource constraints which only increases the difficulty of security testing.
The Mirai botnet~\cite{antonakakis2017mirai} demonstrated the consequences when hundreds of thousands of such devices were conscripted into a DDoS network
that disrupted major Internet infrastructure.
A single vendor ships hundreds of firmware builds across its product line, each with minute adjustments, patches, and changes that make it slightly
different from similar firmware in its cohort.
Manual auditing cannot keep pace with this volume: the global cybersecurity workforce gap reached 4.8 million unfilled positions in
2024~\cite{isc2workforce2024}, and firmware security demands specialized binary analysis skills beyond normal user-land programs.

The dominant automated approach to finding taint-style vulnerabilities in firmware (such as command injection and buffer overflows) is to trace data
flows from attacker-controlled sources to dangerous sinks.
Tools like Karonte~\cite{redini2020karonte}, SaTC~\cite{chen2021satc}, DTaint~\cite{cheng2018dtaint}, and Saluki~\cite{gotovchits2018saluki} all
adopt this approach.
The central limitation of this approach is that taint tracking loses information of how values reach the sink, intact or not.
When an analyst traces a value from a network input to a sink, they reason about what happened to that value along the way.
Integer conversion, truncation, or sanitization are just a few of the factors analysts use to judge when determining if a path is exploitable.
Conventional binary taint tracking does not encode these judgments.
Tools that track only tainted bits produce false positives: they flag paths where the value has been sanitized beyond exploitation \(over tainting\).
Studies of static analysis adoption have repeatedly identified false positives as the primary barrier to practical
use~\cite{johnson2013staticanalysis}.
Existing tools for firmware static analysis attempt to reduce false positives by restricting their analysis scope to achieve tractable running times, analyzing only ``border binaries'' that appear to handle
network input~\cite{chen2021satc, redini2020karonte}.
This filtering discards the majority of the firmware.
On the same dataset, Karonte and SaTC analyze a maximum of 131 of 3,599 binaries.
Vulnerabilities in the remaining binaries go unexamined.

In Chapter~\ref{ch:mango}, I present Operation Mango~\cite{gibbs2024mango}, a static data-flow analysis engine that encodes value-level reasoning into
firmware vulnerability discovery.
\mango's core analysis, \rda, tracks the full transformation history of values through Rich Expressions rather than binary taint bits.
Combined with backward tracing from sinks (Sink-to-Source) and selective callee skipping (Assumed Nonimpact), Mango analyzes every binary in a firmware image.
On the Karonte dataset, Mango covered all 3,599 binaries while averaging 8 minutes per binary, compared to SaTC's 6.5 hours.
On a large-scale corpus of 1,698 firmware images (770,374 binaries), Mango produced 10,834 ranked vulnerability alerts and confirmed 166 exploitable
vulnerabilities in a sampled audit.


\section{Autonomous Cyber Reasoning Systems}

Confirming that a vulnerability is real requires more than any single analysis technique can provide.
Fuzzers produce crashes that need diagnosis.
A segmentation fault does not, by itself, reveal the root cause, the vulnerable input field, or whether
the bug is security-relevant.
Static analyzers flag suspicious patterns without execution context, and studies of their adoption consistently identify false positives as the primary
barrier to practical use~\cite{johnson2013staticanalysis}.
LLM-based scanners identify code that resembles known vulnerability classes but cannot verify their findings dynamically~\cite{thapa2022vuldeepecker}.
Each technique captures a different facet of the problem, alone, each tool produces partial evidence that a human analyst must evaluate, correlate, and
either confirm or discard.

Observational studies of vulnerability researchers show how practitioners close this gap~\cite{votipka2019observational}.
Analysts do not run a single tool and accept its verdict.
They fuzz a target, scan its source code, review crash reports, and cross-reference findings across all three.
When a fuzzer crashes on a function that a static analyzer also flagged, and manual review confirms the root cause, the analyst has built a case through
convergent evidence that no single technique could produce.
This multi-technique coordination, the understanding of how to sequence independent analyses so that their findings reinforce one another, is a form of
practitioner knowledge that current automated systems discard.
Each tool operates in its own silo, producing outputs that a human analyst must manually integrate, and the expertise required for that integration
remains one of the scarcest resources in security~\cite{isc2workforce2024}.

DARPA's Cyber Grand Challenge (CGC)~\cite{darpa2016cgc} in 2016 produced the first generation of autonomous Cyber Reasoning Systems (CRS) that
attempted to replace human integration.
Systems such as Mechanical Phish~\cite{shoshitaishvili2018mechanical} and Mayhem~\cite{cha2012mayhem} demonstrated that autonomous
vulnerability discovery, exploitation, and patching were possible, but only under the constraints of a purpose built operating system (DECREE)
with simplified binaries, input formats, and no library dependencies.
DARPA's AI Cyber Challenge (AIxCC)~\cite{darpa2025aixcc}, a \$29.5 million competition launched in 2024, challenged
this limitation by targeting real-world open-source software: the Linux kernel, SQLite, Jenkins, and similar projects whose scale, language diversity,
and input complexity exceeded what any single technique could handle alone.
AIxCC required systems that find, validate, and patch vulnerabilities with zero human intervention, forcing teams to encode the multi-technique workflow
knowledge that human security teams apply naturally.

In Chapter~\ref{ch:artiphishell}, I present Artiphishell, a fully autonomous CRS designed and led by team Shellphish for AIxCC.
Artiphishell encodes multi-technique workflow knowledge by coordinating coverage-guided fuzzing, grammar-based fuzzing, LLM-powered vulnerability
scanning, statistical and LLM-based root-cause analysis, and AI-driven patching through a multi-agent pipeline.
All agents share state through the Analysis Graph, a central knowledge base that enables the same cross-referencing that human teams perform.
The system builds convergent evidence for a vulnerability, for example: a fuzzer finds a crash in a function that the LLM scanner also flagged, root-cause analysis
localizes the bug, and a patch is generated which passes the validation pipeline without introducing regressions. 


\section{Closing the Loop: Symbolic Execution for Path Feasibility}

Scalable static analyzers for firmware, including \toolname, \satc, \karonte, etc. produce vulnerability candidates, but they cannot prove that any individual candidate is
reachable under conditions that trigger the bug.
The data flow may exist, and the value at the sink may appear exploitable, yet the specific sequence of branches required to reach that sink
could be unsatisfiable given the program's actual constraints.
In firmware, this problem is compounded by cross-binary communication: a value may pass through NVRAM, shared memory, or a socket between processes,
and the analyst must reason about feasibility across multiple binaries with limited visibility into each one.
A human analyst carefully traces control flow, evaluates branch conditions, and judges whether a concrete input could drive the program along the flagged path.
This path feasibility reasoning by humans is expensive, time-consuming, and requires domain expertise.

Symbolic execution~\cite{king1976symbolic} automates this path-based reasoning by treating program inputs as symbolic variables and collecting path constraints that,
when solved, produce concrete triggering inputs.
However, path explosions~\cite{krishnamoorthy2010tackling} make full-program symbolic execution intractable for anything beyond a toy program.
Deeply nested loops, indirect calls through function tables, and library code that the analysis must model or explore significantly hinders progress.
Systems like Driller~\cite{stephens2016driller} and Arbiter~\cite{vadayath2022arbiter} have shown that combining symbolic execution with other
techniques can recover precision that either loses in isolation, but these hybrid approaches operate on low-level intermediate
representations.
Operating at the granularity of assembly and assembly-like representations produces unnecessarily complex SMT constraints, causing executions to stagnate or terminate prematurely.

In Chapter~\ref{ch:proposed}, I propose a symbolic execution engine that operates on angr's AIL (angr Intermediate Language), a higher-level IR that
carries recovered type information and structured control flow in the form of a near-decompilation representation.
AIL-level symbolic execution simplifies path feasibility reasoning by dropping Single Static Assignment (SSA), allowing
constraints from recovered types, structured loop and conditional handling to operate on the variable level of source code.
The engine consumes static analysis results and proves with high accuracy that a vulnerability candidate is reachable under conditions that trigger the bug.
Leveraging the inter-program communication networks of firmware, we can greatly simplify the path constraints by treating inter-program communication as a form of symbolic input,
allowing us to bypass complex path conditions that would otherwise require deep symbolic execution.
For conditions and sub-routines that still require deep symbolic execution, we substitute symbolic execution with LLM generated subroutines.
This allows us to greatly increase the speed and fidelity of our analysis.
